\documentclass[11pt]{article}
\usepackage[T1]{fontenc}
\usepackage[utf8]{inputenc}
\usepackage[brazil]{babel}
\usepackage{amsmath,amssymb,booktabs,longtable,array,geometry,graphicx,hyperref,microtype}
\geometry{margin=25mm}
\hypersetup{colorlinks=true,linkcolor=blue,urlcolor=blue,citecolor=blue}
\newcommand{\vet}[1]{\boldsymbol{#1}}
\newcommand{\codigo}[1]{\texttt{#1}}
\newcommand{\figpendente}[2]{\par\medskip\noindent\fbox{\parbox{0.94\textwidth}{\textbf{FIGURA A PREPARAR:} #1\\\textbf{Base sugerida:} #2}}\par\medskip}

\title{Nota Técnica do FreeCAD Chain Drive Workbench\\
Conversão de catálogos comerciais em modelos CAD e solução geométrica discreta de transmissões por corrente}
\author{Douglas D. Schons}
\date{Versão de trabalho para revisão --- 31 de agosto de 2026}

\begin{document}
\maketitle

\begin{abstract}
O \emph{FreeCAD Chain Drive Workbench} foi desenvolvido para resolver dois
problemas recorrentes no projeto mecânico de transmissões por corrente. O
primeiro é transformar informações dispersas em catálogos comerciais e normas
técnicas em modelos CAD tridimensionais utilizáveis diretamente no FreeCAD.
Essa função permite selecionar uma corrente ASA e rodas dentadas comerciais e
gerar a corrente, uma roda dentada isolada ou o conjunto completo formado por
duas rodas dentadas e uma corrente. O segundo problema é a incompatibilidade
entre uma distância entre centros inicialmente desejada e o comprimento
discreto da corrente, que somente pode possuir um número inteiro de passos. O
Workbench calcula o número de elos, identifica a necessidade de elo de redução
e corrige a distância entre centros para produzir um polígono fechado de
centros de roletes com passo constante. Esta nota apresenta o fluxo completo
do catálogo ao CAD, as condições de contorno comerciais, o equacionamento
geométrico, a geração paramétrica dos elos e das rodas dentadas, a integração
com o FreeCAD e a metodologia de validação.
\end{abstract}

\noindent\textbf{Palavras-chave:} corrente de rolos; roda dentada; sprocket;
distância entre centros; modelo discreto; catálogo comercial; FreeCAD; CAD
paramétrico; ASA; ANSI.

\tableofcontents

\section{O problema de engenharia resolvido pelo Workbench}

No desenvolvimento de uma transmissão por corrente, o projetista normalmente
precisa realizar duas tarefas que, embora relacionadas, costumam ser executadas
separadamente.

A primeira tarefa é selecionar componentes comerciais. Catálogos de
fabricantes apresentam séries ASA, passos, diâmetros de rolete, larguras,
números de dentes disponíveis, diâmetros externos, cubos e limites de furação.
Essas informações são adequadas para consulta, mas não constituem diretamente
um conjunto CAD. Para utilizá-las em um projeto mecânico, o projetista ainda
precisa interpretar tabelas, reconstruir geometrias e organizar as peças no
ambiente tridimensional.

A segunda tarefa é fechar geometricamente o acionamento. A distância entre
centros desejada pode ser qualquer valor real, mas a corrente somente pode ser
montada com um número inteiro de intervalos de passo. Quando a geometria é
tratada apenas por comprimentos contínuos de arcos e tangentes, o resultado é
uma boa estimativa inicial, porém não assegura que todos os centros de roletes
estejam separados exatamente por um passo e que o último elo encontre o
primeiro sem deformação artificial.

O Workbench reúne essas duas tarefas em um único fluxo:
\begin{enumerate}
\item consulta uma base de dados derivada de catálogos comerciais;
\item limita as escolhas às correntes e rodas dentadas disponíveis nessa base;
\item calcula a geometria discreta do acionamento;
\item determina o número de elos e a distância entre centros corrigida;
\item gera corpos CAD paramétricos e os organiza em um documento FreeCAD.
\end{enumerate}

Portanto, o programa funciona simultaneamente como um catálogo digital
orientado a CAD e como uma ferramenta de resolução geométrica do acionamento.

\figpendente{Visão geral do problema: catálogo comercial à esquerda, cálculo discreto da distância entre centros ao centro e conjunto CAD no FreeCAD à direita.}{Criar uma figura nova combinando uma tabela simplificada do catálogo ENCO, o diagrama \codigo{polygonal\_layout\_asa80\_11\_20.png} do repositório \emph{chain-drive-geometry-calculator} e uma captura do Workbench atual.}

\section{Resultados fornecidos ao usuário}

O Workbench oferece três saídas principais:
\begin{itemize}
\item \textbf{Roda dentada isolada:} gera uma roda dentada comercial da série
ASA selecionada, com número de dentes disponível no catálogo e furo dentro dos
limites cadastrados;
\item \textbf{Corrente isolada:} gera uma corrente fechada, formada por elos
internos, externos e, quando necessário, um elo de redução, já posicionados de
acordo com a solução do entrecentros;
\item \textbf{Conjunto completo:} gera duas rodas dentadas e a corrente em um
único documento, com a roda menor em (X=0), a roda maior na distância
corrigida e as peças faseadas em relação aos centros dos roletes.
\end{itemize}

Antes da geração tridimensional, o usuário recebe o número de elos, a indicação
de elo de redução, a distância desejada, a distância corrigida, a correção
necessária, o comprimento discreto (Np), o resíduo de fechamento e o erro
máximo de passo.

\figpendente{As três opções do Workbench: roda dentada, corrente e conjunto completo.}{Capturas novas dos três comandos registrados por \codigo{chain\_drive\_commands.py}; usar a mesma série ASA e o mesmo padrão visual.}

\section{Condição de contorno: componentes ASA comerciais}

O escopo atual não é um gerador universal de correntes e rodas dentadas. A
condição de contorno do projeto é a disponibilidade comercial registrada nos
catálogos adotados. Primeiro é selecionada a série ASA; em seguida, o programa
apresenta somente números de dentes e dimensões para os quais existe um registro
válido.

Essa barreira de entrada é importante. Uma combinação pode ser matematicamente
possível sem corresponder a uma peça comercial considerada no projeto. O
Workbench não interpola silenciosamente dimensões ausentes e não cria uma roda
dentada comercial inexistente. A matemática opera dentro do domínio definido
pelas tabelas instaladas.

As dimensões teóricas de passo e do vão entre dentes são combinadas com
dimensões comerciais de corpo. Assim, deve-se distinguir:
\begin{itemize}
\item \textbf{dimensões da corrente}, associadas à série ASA;
\item \textbf{geometria teórica de engrenamento}, associada ao passo, ao
diâmetro do rolete e ao número de dentes;
\item \textbf{dimensões comerciais da roda}, como diâmetro externo, cubo,
comprimento e limites de furação.
\end{itemize}

\section{Fluxo do catálogo ao modelo CAD}

O fluxo de dados pode ser resumido por
\[
\text{catálogo}\rightarrow\text{registro validado}\rightarrow
\text{cálculo geométrico}\rightarrow\text{parâmetros CAD}\rightarrow
\text{corpos e conjunto FreeCAD}.
\]

O arquivo \codigo{data/chains/enco\_asa\_chains.csv} contém dimensões das
correntes. Os arquivos em \codigo{data/sprockets/enco/} contêm as rodas
dentadas por série. Os módulos \codigo{core/catalog.py} e
\codigo{core/sprocket\_catalog.py} normalizam separadores decimais, verificam
campos e disponibilizam apenas registros válidos. O módulo
\codigo{application.py} reúne o catálogo, o solucionador e os geradores CAD.

\figpendente{Fluxo de informação e módulos: arquivos de catálogo, núcleo matemático, geradores CAD, interface e documento FreeCAD.}{Desenhar a partir de \codigo{application.py} e \codigo{docs/ARCHITECTURE\_INVENTORY.md}; não requer figura externa.}

\section{Dados dimensionais da corrente}

Para cada série ASA, os parâmetros geométricos efetivamente consumidos são:
\begin{table}[ht]
\centering
\caption{Dimensões da corrente utilizadas na geração CAD.}
\begin{tabular}{lll}\toprule
Símbolo & Campo & Significado\\\midrule
(p)&\codigo{pitch\_P\_mm}&passo\\
(E)&\codigo{inner\_width\_E\_mm}&largura interna livre\\
(D_R)&\codigo{roller\_diameter\_R\_mm}&diâmetro do rolete\\
(H)&\codigo{plate\_height\_H\_mm}&altura da placa\\
(D_G)&\codigo{pin\_diameter\_G\_mm}&diâmetro do pino\\
(L)&\codigo{overall\_width\_L\_mm}&largura total\\
(T_p)&\codigo{plate\_thickness\_T\_mm}&espessura da placa\\\bottomrule
\end{tabular}
\end{table}

Para a corrente ASA80 usada no caso canônico, (p=25{,}40) mm,
(E=15{,}75) mm, (D_R=15{,}88) mm, (H=24{,}00) mm,
(D_G=7{,}92) mm, (L=32{,}70) mm e (T_p=3{,}25) mm.

\figpendente{Corrente de rolos cotada com (p,E,D_R,H,D_G,L,T_p).}{Redesenhar a anatomia mostrada no manual iwis, p. 9, ou no guia Tsubaki, substituindo a nomenclatura pela adotada nesta nota.}

\section{Dados comerciais da roda dentada}

Cada registro de roda dentada contém número de dentes (z), diâmetro de passo
de catálogo para conferência, diâmetro externo, diâmetro do cubo, comprimento
total, furo piloto e furo máximo. O furo escolhido deve satisfazer
\begin{equation}
D_{piloto}\leq D_f\leq D_{máximo}.
\end{equation}

Esses valores definem o corpo comercial. O perfil do vão entre dentes é
calculado separadamente e subtraído desse corpo. Essa separação permite que uma
mesma formulação teórica seja aplicada às opções comerciais cadastradas sem
confundir diâmetro de passo com diâmetro do bruto.

\figpendente{Desenho comercial da roda dentada com diâmetro externo, cubo, largura da flange, comprimento e furos.}{Localizar a página específica da série no \codigo{catalogo\_enco.pdf} e redesenhar; não inserir a página escaneada.}

\section{Entradas do cálculo do acionamento}

As entradas principais são
\begin{equation}
p,\qquad z_1,\qquad z_2,\qquad C_d,
\end{equation}
onde (p) é obtido do catálogo da corrente, (z_1) e (z_2) são números de
dentes comercialmente disponíveis, e (C_d) é a distância entre centros
desejada. O número total de elos (N) não é imposto inicialmente. Ele é
estimado pela geometria contínua e testado como inteiro no modelo discreto.

O usuário também pode exigir um número par de elos. Quando essa restrição não é
ativada, uma solução com (N) ímpar é permitida e requer um elo de redução
(\emph{OffsetLink}).

\section{Geometria de passo das rodas dentadas}

Os centros dos roletes assentados formam um polígono regular. Para uma roda
com (z) dentes,
\begin{equation}
r=\frac{p}{2\sin(\pi/z)},\qquad D_p=2r=\frac{p}{\sin(\pi/z)}.
\label{eq:raio_passo}
\end{equation}
Logo,
\begin{equation}
r_1=\frac{p}{2\sin(\pi/z_1)},\qquad
r_2=\frac{p}{2\sin(\pi/z_2)}.
\end{equation}
O incremento angular entre posições consecutivas é
\begin{equation}
\delta_1=\frac{2\pi}{z_1},\qquad
\delta_2=\frac{2\pi}{z_2}.
\end{equation}

\figpendente{Polígono de passo com (p,r,D_p,delta).}{Adaptar \codigo{ASAChainDriveApp/docs/figures/sprocket\_fig\_02\_pitch\_geometry.pdf}; comparar com a figura de geometria do repositório antigo.}

\section{Estimativa contínua inicial}

Definindo (Delta r=r_2-r_1), para uma distância (C),
\begin{equation}
\alpha(C)=\arcsin\left(\frac{\Delta r}{C}\right),\qquad
T(C)=\sqrt{C^2-\Delta r^2}.
\end{equation}
Os ângulos contínuos de contato são
\begin{equation}
\beta_1(C)=\pi-2\alpha(C),\qquad
\beta_2(C)=\pi+2\alpha(C),
\end{equation}
e o comprimento de referência é
\begin{equation}
L_{cont}(C)=r_1\beta_1(C)+r_2\beta_2(C)+2T(C).
\label{eq:comprimento_continuo}
\end{equation}
Assim,
\begin{equation}
N_0=\operatorname{round}\left(\frac{L_{cont}(C_d)}{p}\right).
\end{equation}

Essa formulação é somente uma estimativa. Sobre os arcos, o comprimento de arco
entre dois pontos é diferente do comprimento da corda rígida que representa um
elo. A solução física final deve impor o passo entre centros de roletes.

\figpendente{Geometria contínua com raios, tangentes, ângulos de contato e distância entre centros.}{Adaptar a figura conceitual da formulação em \codigo{complete\_mathematical\_formulation\_en.tex}; redesenhar com a nomenclatura portuguesa.}

\section{Fundamento discreto: roletes e intervalos de passo}

O projeto anterior \emph{chain-drive-geometry-calculator} estabeleceu uma
distinção essencial: (M) roletes em contato com uma roda definem (M-1)
intervalos de passo no arco. Definindo (M_1,M_2) como números de roletes em
contato e (S,R) como intervalos nos ramos superior e inferior,
\begin{equation}
(M_1-1)+S+(M_2-1)+R=N,
\label{eq:topologia}
\end{equation}
ou (M_1+M_2+S+R=N+2). Essa equação descreve a topologia discreta e continua
útil para explicar por que o número de elos não pode ser tratado como um
comprimento contínuo.

Para cada roda,
\begin{equation}q=\frac{M-1}{2}.\end{equation}
Se (M) é ímpar, há um rolete no eixo de simetria; se (M) é par, o eixo passa
entre dois roletes. Essa diferença é visível nas figuras do projeto anterior e
deve ser mantida nas novas ilustrações.

\figpendente{Comparação entre contato com (M) ímpar e (M) par, mostrando roletes e intervalos.}{Adaptar \codigo{polygonal\_layout\_asa80\_11\_20.png} e \codigo{polygonal\_layout\_asa80\_19\_20.png} do repositório \emph{chain-drive-geometry-calculator}.}

\section{Formulação topológica direta do projeto anterior}

No sistema (O_1=(0,0)), (O_2=(C,0)), os extremos simétricos de contato da
roda menor podem ser escritos como
\begin{align}
A_t&=O_1+r_1[\cos(\pi-q_1\delta_1),\ \sin(\pi-q_1\delta_1)]^T,\\
A_b&=O_1+r_1[\cos(\pi+q_1\delta_1),\ \sin(\pi+q_1\delta_1)]^T,
\end{align}
com (q_1=(M_1-1)/2). Para a roda maior,
\begin{align}
B_t(C)&=O_2+r_2[\cos(q_2\delta_2),\ \sin(q_2\delta_2)]^T,\\
B_b(C)&=O_2+r_2[\cos(-q_2\delta_2),\ \sin(-q_2\delta_2)]^T,
\end{align}
com (q_2=(M_2-1)/2).

Os ramos retos satisfazem
\begin{equation}
\|B_t-A_t\|=Sp,\qquad \|A_b-B_b\|=Rp.
\end{equation}
Para uma topologia válida, a distância entre centros pode ser obtida diretamente:
\begin{equation}
C=A_{t,x}-r_2\cos(q_2\delta_2)+
\sqrt{(Sp)^2-[r_2\sin(q_2\delta_2)-A_{t,y}]^2}.
\label{eq:centro_topologico}
\end{equation}

Essa formulação explica matematicamente a seleção de topologias e permanece
como base documental do problema. O Workbench atual usa uma implementação
mais geral por caminhada de cordas sobre o caminho de passo, descrita a seguir.
As duas formulações não devem ser apresentadas como se fossem o mesmo algoritmo:
a primeira explicita (M_1,S,M_2,R); a segunda fecha diretamente uma sequência
de (N) passos rígidos.

\section{Modelo discreto implementado no Workbench}

Seja (\vet p_i) o centro do rolete (i). A condição fundamental é
\begin{equation}
\|\vet p_{i+1}-\vet p_i\|=p,qquad i=0,\ldots,N-1.
\label{eq:passo_rigido}
\end{equation}
O caminho geométrico (et\gamma(s)) conserva os arcos de passo e as tangentes
externas. Partindo de (s_i), o próximo rolete é a primeira solução à frente de
\begin{equation}
f_i(\Delta s)=\|\vet\gamma(s_i+\Delta s)-\vet\gamma(s_i)\|-p=0.
\label{eq:proximo_rolete}
\end{equation}

O código procura uma mudança de sinal em incrementos (p/50), limitada a
(5p), e aplica bisseção. Dessa maneira, os centros permanecem sobre a linha
de passo, mas a distância entre centros consecutivos é a corda rígida (p), e
não o comprimento de arco.

\figpendente{Ampliação da transição arco--reta, mostrando incremento de caminho (Delta s) e corda rígida (p).}{Gerar a partir de \codigo{core/discrete\_solver.py}; usar a estética dos gráficos \codigo{plot\_1.png}/\codigo{plot\_2.png} do projeto anterior.}

\section{Seleção do número inteiro de elos}

A partir de (N_0), a implementação testa, por padrão,
\begin{equation}
N\in[N_0-3,N_0+3],\qquad N\geq4.
\end{equation}
Se a opção de corrente par estiver ativa, candidatos ímpares são removidos.
Cada candidato é resolvido para sua própria distância (C_N). A solução
selecionada minimiza
\begin{equation}|C_N-C_d|.\end{equation}
Se (N) for ímpar, \codigo{requires\_offset\_link=True} e o último intervalo
recebe um OffsetLink.

\section{Fechamento e correção da distância entre centros}

Após caminhar (N) cordas, seja (s_N(C)) a coordenada não reduzida módulo o
comprimento total do caminho (L_\Gamma(C)). A função escalar implementada é
\begin{equation}
R_s(C;N)=s_N(C)-L_\Gamma(C).
\label{eq:residuo_fechamento}
\end{equation}
O sinal informa se a caminhada ultrapassou ou não completou uma volta. O
programa expande um intervalo em torno de (C_d) e aplica bisseção até
(R_s\approx0). Então, pela periodicidade,
\begin{equation}
\vet p_N\approx\vet p_0.
\end{equation}

É importante distinguir o resíduo assinado (R_s), armazenado atualmente como
\codigo{closure\_residual\_mm}, do erro euclidiano
(e_c=\|\vet p_N-\vet p_0\|). O primeiro é a função usada para resolver (C);
o segundo é uma interpretação geométrica do fechamento.

O erro máximo de passo é
\begin{equation}
e_{max}=\max_i\left|\|\vet p_{i+1}-\vet p_i\|-p\right|.
\end{equation}

\figpendente{Fluxograma: estimativa contínua, candidatos (N), caminhada de passos, raiz em (C), seleção por menor correção.}{Recriar a partir de \codigo{core/discrete\_solver.py}, mantendo a sequência conceitual da documentação antiga.}

\section{Sistema de coordenadas e poses CAD}

O plano do solucionador mapeia-se no FreeCAD por
\begin{equation}(x,y)_{solver}\mapsto(X,Y,Z)_{CAD}=(x,0,y).\end{equation}
Um ângulo positivo do solucionador é uma rotação em torno de (-Y). Cada
template possui
\begin{equation}
JointA=(0,0,0),\qquad JointB=(p,0,0).
\end{equation}
Para o intervalo (i),
\begin{equation}
\theta_i=\operatorname{atan2}(y_{i+1}-y_i,x_{i+1}-x_i).
\end{equation}
A translação é ((x_i,0,y_i)) e a rotação é (	heta_i) em torno de (-Y).

\figpendente{Mapeamento do plano matemático para o plano (XZ) do FreeCAD e convenção JointA/JointB.}{Criar com base nos comentários de \codigo{cad/chain\_generator.py}.}

\section{Geração do elo interno}

O InnerLink contém duas placas internas e dois roletes vazados. As placas são
perfis simétricos de quatro arcos extrudidos ao longo de (Y). O raio da
extremidade é (H/2); a semialtura de pescoço implementada é
(h_n=0{,}875H/2), e o raio de alívio é
\begin{equation}
r_q=\frac{(H/2)^2-(p/2)^2-h_n^2}{2(h_n-H/2)}.
\end{equation}
As faces internas ficam em (Y=\pm E/2). Os roletes possuem diâmetro externo
(D_R), furo (D_G) e comprimento (E).

\figpendente{InnerLink explodido, componentes, eixos e interfaces.}{Captura nova de \codigo{ASA80\_InnerLink.FCStd} ou template atual.}

\section{Geração do elo externo}

O OuterLink utiliza o mesmo perfil de placa. As placas externas são deslocadas
para fora do envelope do elo interno, e dois pinos maciços de diâmetro (D_G)
atravessam a largura total (L) nas posições (X=0) e (X=p). Detalhes de
rebitamento, ajustes e estampagem específicos de fabricante não são modelados.

\figpendente{OuterLink explodido com placas e pinos.}{Captura nova de \codigo{ASA80\_OuterLink.FCStd} ou template atual.}

\section{Geração do elo de redução}

O OffsetLink fecha cadeias com (N) ímpar. Seu template canônico possui pino em
JointA e rolete em JointB. As placas deslocadas são construídas por loft entre
seções retangulares, com transição lateral suave. O sólido inclui placas,
rolete, bucha e pino de conexão.

O último intervalo do solucionador percorre rolete para pino, sentido oposto ao
template. Na montagem, a origem é movida para o final do intervalo e a rotação
recebe (180^\circ). Essa operação preserva os dois centros calculados sem
alterar a geometria validada do OffsetLink.

\figpendente{OffsetLink em vista lateral/isométrica e reversão no último intervalo.}{Usar \codigo{ASAChainDriveApp/output/validation/ASA80\_OffsetLink\_Standalone\_*.FCStd}.}

\section{Montagem eficiente da corrente}

O documento cria apenas três corpos mestres: InnerLink, OuterLink e OffsetLink.
Cada elo da corrente é uma instância \codigo{App::Link} de um desses templates.
Isso reduz a duplicação de representações geométricas e permite montar dezenas
ou centenas de elos mantendo propriedades individuais de índice, tipo, ângulo e
mapeamento de interface.

\section{Geometria de assentamento do rolete}

O diâmetro e o raio de assentamento implementados são
\begin{equation}
D_s=1{,}005D_R+0{,}003(25{,}4\ \mathrm{mm}),\qquad R_s=D_s/2.
\end{equation}
No sistema local do vão, o ponto de passo é (a=(0,r)), e o ponto mais profundo
é (L=(0,r-R_s)). Essas expressões devem ser conferidas com a edição controlada
da ASME B29.1 antes da documentação final de fabricação.

\figpendente{Círculo de assentamento, ponto (a), raio (R_s), ponto (L) e rolete.}{Adaptar \codigo{ASAChainDriveApp/docs/figures/sprocket\_fig\_03\_seating\_curve.pdf} e conferir com GEARS pp. 2--4.}

\section{Construção teórica do vão entre dentes}

Com (n=z) e (phi=\pi/n), o código utiliza
\begin{align}
A&=35^\circ+60^\circ/n,&B&=18^\circ-56^\circ/n,\\
ac&=0{,}8D_R,&M&=ac\cos A,&T&=ac\sin A,\\
E_t&=1{,}3025D_R+0{,}0015(25{,}4\ \mathrm{mm}),\\
ab&=1{,}4D_R,&W&=ab\cos\phi,&V&=ab\sin\phi,\\
yz&=D_R[1{,}4\sin(17^\circ-64^\circ/n)-0{,}8\sin(18^\circ-56^\circ/n)],\\
F&=D_R[0{,}8\cos(18^\circ-56^\circ/n)+1{,}4\cos(17^\circ-64^\circ/n)-1{,}3025]
-0{,}0015(25{,}4\ \mathrm{mm}).
\end{align}

São construídos os pontos (a,b,c,x,y,z,L). O ponto (x) é a interseção mais
distante da reta (cx) com o círculo de assentamento; o centro do arco (xy)
é a solução de raio (E_t) mais próxima de (c); entre as duas possibilidades
para (z), escolhe-se a que satisfaz melhor o círculo de raio (F) centrado em
(b). O perfil passa pela curva de topo (F), segmento (zy), arco (E_t) e
arco de assentamento até (L), sendo espelhado em relação ao eixo radial.

\figpendente{Construção completa (a,b,c,x,y,z,L), ângulos, raios e eixo de espelhamento.}{Redesenhar a Fig. 1 e as etapas das pp. 3--12 de \codigo{design\_draw\_sprocket\_5.pdf}; usar \codigo{sprocket\_fig\_01\_reference\_system.pdf} como base secundária.}

\section{Geração paramétrica da roda dentada}

O corpo axial combina diâmetro externo, cubo, comprimento e furo comerciais
com largura de flange e chanfro da seção dentada. O perfil axial é revolucionado
em torno de (Y). Um cortador de vão é extrudido além das faces, copiado (z)
vezes e rotacionado em passos (360^\circ/z). A subtração booleana resulta em
um único sólido válido. Rodas já geradas podem ser reutilizadas por uma
biblioteca local associada à série, dentes, furo e versão do perfil.

\section{Geração do conjunto completo}

A roda menor é posicionada em (X=0), e a maior em (X=C_N). Para cada roda,
o programa identifica o centro de rolete mais próximo do respectivo raio de
passo e usa seu ângulo polar como fase de montagem. Essa regra alinha um vão com
um rolete no estado estático. Não representa uma simulação dinâmica de
engrenamento.

\figpendente{Conjunto canônico com rodas, corrente, distância desejada/corrigida e tipos de elo.}{Captura nova do Workbench para ASA80, 11/20 dentes e (C_d=400) mm.}

\section{Integração com o FreeCAD}

\codigo{InitGui.py} registra a bancada e três comandos. A interface unificada
carrega as opções válidas, impede combinações fora do catálogo, calcula os
resultados e invalida cálculos antigos quando uma entrada muda.
\codigo{application.py} constitui a API de alto nível. O núcleo matemático e os
leitores de catálogo não importam FreeCAD; os módulos CAD são carregados somente
durante a geração.

\section{Caso canônico de validação}

Para ASA80, (p=25{,}40) mm, (z_1=11), (z_2=20) e (C_d=400) mm, a
execução atual produz
\begin{align}
N&=47,&\text{OffsetLink}&=\text{sim},\\
C_N&=398{,}338658800570\ \mathrm{mm},&
\Delta C&=-1{,}661341199430\ \mathrm{mm},\\
R_s&=2{,}1100\times10^{-10}\ \mathrm{mm},&
e_{max}&=2{,}1042\times10^{-10}\ \mathrm{mm}.
\end{align}

Os resíduos são métricas numéricas do algoritmo, não tolerâncias físicas. O
caso deve ser complementado por medição independente entre eixos e verificação
do passo no FreeCAD.

\figpendente{Comparação lado a lado entre polígono calculado e medição CAD do caso canônico.}{Adaptar o formato dos pares \codigo{plot\_1.png}/\codigo{cd\_1.png} e \codigo{plot\_2.png}/\codigo{cd\_2.png} do projeto antigo, substituindo pelo caso atual 11/20/400.}

\section{Metodologia de validação}

Cada caso deve registrar versão do código, versão do FreeCAD, registros de
catálogo, (p,z_1,z_2,C_d,N,C_N,\Delta C), paridade, (R_s), (e_{max}),
validade dos sólidos e medidas independentes no CAD. A validação é dividida em:
\begin{enumerate}
\item invariantes matemáticos e testes puros do solucionador;
\item validação de leitura e seleção dos catálogos;
\item geração headless e validade OpenCascade dos sólidos;
\item teste dos três fluxos da interface;
\item medição independente dos modelos gerados.
\end{enumerate}

Uma matriz numérica já cobre séries ASA35 a ASA160, mas a validação CAD completa
está concentrada na ASA80. Esses níveis não devem ser confundidos.

\section{Limitações}

O Workbench atual considera duas rodas, eixos paralelos, centros alinhados,
transmissão aberta, corrente simples e opções presentes no catálogo. Não inclui
esticador, roda intermediária, flecha, elasticidade, desgaste, folga, dinâmica,
forças de contato, potência, tensão, fadiga, lubrificação ou proteção de máquina.
Os corpos CAD são representações de projeto; não incorporam todos os detalhes
de fabricação, ajustes, materiais, tratamentos ou tolerâncias.

O OffsetLink contém proporções funcionais validadas principalmente na ASA80.
As constantes do perfil dentado e as larguras de flange ainda devem ser
conferidas contra uma edição controlada da ASME B29.1. As tabelas derivadas do
ENCO ainda precisam de rastreabilidade página a página.

\section{Trabalhos futuros}

As próximas etapas são: adaptar as figuras do projeto matemático anterior;
medir vários casos no FreeCAD; validar OffsetLink em outras séries; completar a
rastreabilidade do catálogo; conferir as constantes normativas; ampliar testes
de rodas nos extremos de dentes; e somente depois preparar README, publicação e
distribuição. Extensões futuras podem tratar correntes múltiplas, cubos
alternativos, esticadores e dimensionamento de potência, mantendo essas análises
separadas do núcleo geométrico já validado.

\section{Conclusão}

O FreeCAD Chain Drive Workbench transforma dados comerciais de correntes e
rodas dentadas ASA em modelos CAD organizados e resolve a incompatibilidade
entre uma distância entre centros arbitrária e um número inteiro de passos. A
estimativa contínua orienta a busca; a condição final é discreta,
\(|\vet p_{i+1}-\vet p_i|=p), com fechamento após (N) intervalos. Essa
combinação de catálogo, matemática e geração paramétrica reduz consultas
manuais, reconstrução repetitiva de peças e tentativa e erro na montagem CAD.

\section{Referências}
\begin{enumerate}
\item ASME International. \emph{ASME B29.1 --- Precision Power Transmission Roller Chains, Attachments, and Sprockets}. Edição controlada aplicável.
\item ENCO. Catálogo comercial de correntes e rodas dentadas usado na derivação das tabelas instaladas.
\item GEARS Educational Systems. \emph{Sprocket Tooth Design Formulas and Procedure for Drawing a Sprocket}, 15 p., 2010.
\item iwis antriebssysteme GmbH \& Co. KG. \emph{Handbook for Chain Engineering: Design and Construction / Examples of Calculation}.
\item Tsubaki. \emph{Guia de Produtos}, cópia local em português.
\item ISO 606. \emph{Short-pitch transmission precision roller and bush chains, attachments and associated chain sprockets}.
\item Schons, D. D. \emph{Chain Drive Geometry Calculator}. Repositório anterior: \url{https://github.com/douglasdschons/chain-drive-geometry-calculator}.
\item Schons, D. D. \emph{FreeCAD Chain Drive Workbench}. Código-fonte da versão de desenvolvimento 0.1.0.
\item FreeCAD Project Association. Documentação do FreeCAD, Part/OpenCascade, placements e \codigo{App::Link}.
\end{enumerate}

\end{document}
